%AttackName: Primal-Dual Proximal Gradient Descent Attack (PDPGD)

%Input:Image x, label y, initial perturbation r(0), total number of iterations T, learning rates θr and θλ.
%Output:Adversarial perturbation r that minimally perturbs x to cause misclassification by the model.
%Formula:r^{(k+1)} = \Pi^x_C \left( \text{prox}_{\lambda \|\cdot\|,\theta_r} \left( r^{(k)} - \theta_r \nabla_r^{(k)} \right) \right) \newline \lambda^{(k+1)} = \Pi_\Lambda \left( \lambda^{(k)} + \theta_\lambda \nabla_\lambda^{(k)} \right)
%Explanation:The Primal-Dual Proximal Gradient Descent Attack (PDPGD) is a white-box adversarial attack that directly solves the original non-convex constrained norm minimization problem for generating adversarial examples. It simultaneously updates the primal variable (the adversarial perturbation r) and the dual variable (the penalty λ for misclassification constraint) using a primal-dual optimization scheme. For non-smooth norms (such as l∞, l1, l0), it uses a proximal gradient step for r, allowing efficient minimization even when the norm is non-differentiable. The attack iteratively updates r using the proximal operator of the chosen norm and updates λ via projected gradient ascent, ensuring the misclassification constraint is enforced. This approach finds minimal-norm adversarial perturbations more efficiently and accurately than previous methods.